% Volume VI: Machine Learning Mathematics
% Continuity note for Chapter 36.
% Previous dependencies: Chapter 15 provides conditional expectation and Bayes decisions; Chapter 27 provides stochastic processes, transition kernels, Markov chains, and filtering intuition; Chapter 30 provides controlled systems, dynamic programming, Bellman equations, and optimal control; Chapter 31 provides risk minimization and losses.
% New durable concepts: Markov decision processes; states, actions, rewards, transition kernels, policies, trajectories, returns, discounting, value functions, Q-functions, Bellman expectation and optimality equations, Bellman backups, temporal-difference errors, Q-learning, policy gradients, score-function identity, actor-critic structure, exploration, model-based versus model-free learning, reward prediction errors, and neural decision dynamics.
% Forward hooks: Chapter 37 treats policies and value networks as parameterized function classes; Chapter 38 differentiates policy and value objectives; Chapter 39 studies optimization dynamics; Chapter 45 links plasticity rules to learning; Chapter 50 connects decisions to Bayesian inference and predictive coding.
% Notation commitments: States are s in S, actions are a in A or A(s), rewards are r_t or R_{t+1}, transition kernel is p(s',r|s,a) or P(s'|s,a), policies are pi(a|s), discount factor is gamma in [0,1), returns are G_t, value functions are V^pi and Q^pi, optimal values are V^* and Q^*, TD error is delta_t.
\section{Chapter 36: Reinforcement Learning}
\label{ch:reinforcement-learning}

Chapters 30 and 35 studied two different kinds of structure. Control theory asks how actions steer a dynamical system. Unsupervised learning asks how data reveal states, geometry, and representations. Reinforcement learning combines these ideas with a third ingredient: reward. An agent observes a state, takes an action, changes the future state distribution, receives reward, and must learn from delayed consequences.

The central difficulty is credit assignment through time. An action may look bad immediately but make later reward possible. Another action may give short-term reward while damaging future value. Reinforcement learning turns this into mathematics through Markov decision processes, value functions, Bellman equations, temporal-difference learning, and policy optimization.

\paragraph{Chapter problem.}
How can an agent learn policies for stochastic sequential decision problems when actions change future states and rewards?

\paragraph{Section map.}
Section 36.1 defines Markov decision processes, policies, trajectories, rewards, and returns. Section 36.2 derives value functions, Bellman equations, Bellman backups, and temporal-difference learning. Section 36.3 develops policy optimization, policy gradients, exploration, and actor-critic structure. Section 36.4 connects RL mathematics to neuroscience through reward prediction errors, action selection, model-based control, and neural decision dynamics.

\subsection{36.1 Markov decision processes}

\paragraph{Guiding question.}
What is the minimal mathematical model of a stochastic environment in which an agent's actions affect future observations and rewards?

\paragraph{Beginner-facing intuition.}
A Markov decision process, or MDP, is a Markov chain with choices. In a Markov chain, the current state determines the distribution of the next state. In an MDP, the agent first chooses an action, and the pair \((\text{state},\text{action})\) determines the distribution of the next state and reward.

The Markov assumption is a modeling choice: the state must contain enough information that the past is irrelevant once the current state and action are known. If the state omits important hidden variables, the process may only appear Markov after state estimation or memory is added.

\paragraph{Formal definitions.}
A finite discounted MDP consists of:
\[
S \quad \text{states},\qquad A(s) \quad \text{available actions at state }s,
\]
a transition and reward kernel
\[
p(s',r\mid s,a)
=P(S_{t+1}=s',R_{t+1}=r\mid S_t=s,A_t=a),
\]
and a discount factor
\[
\gamma\in[0,1).
\]
A stochastic policy is a conditional distribution
\[
\pi(a\mid s)=P(A_t=a\mid S_t=s).
\]
Given a policy, the state-action sequence becomes a stochastic process. The discounted return from time \(t\) is
\[
\boxed{
G_t=R_{t+1}+\gamma R_{t+2}+\gamma^2R_{t+3}+\cdots
=\sum_{k=0}^{\infty}\gamma^k R_{t+k+1}.
}
\]
The discount makes future rewards count less and ensures finite values when rewards are bounded.

\paragraph{Core mechanism: trajectory probabilities factor locally.}
For an initial distribution \(\rho(s_0)\), a policy \(\pi\), and transition kernel \(P(s'\mid s,a)\), the probability of a state-action trajectory
\[
\tau=(s_0,a_0,s_1,a_1,\ldots,s_T)
\]
is
\[
\boxed{
P_\pi(\tau)
=\rho(s_0)\prod_{t=0}^{T-1}\pi(a_t\mid s_t)P(s_{t+1}\mid s_t,a_t).
}
\]
The factorization follows from the chain rule for probability plus the Markov assumption. The past affects the next state only through \(s_t\) and \(a_t\), and the policy samples \(a_t\) from \(s_t\). This local factorization is why dynamic programming and RL algorithms can update one state or transition at a time.

\paragraph{Concrete example.}
Let \(S=\{\mathrm{low},\mathrm{high}\}\) and actions \(A=\{\mathrm{wait},\mathrm{work}\}\). From \(\mathrm{low}\), action \(\mathrm{wait}\) moves to \(\mathrm{high}\) with probability \(0.8\) and reward \(0\). Action \(\mathrm{work}\) gives reward \(1\) but stays \(\mathrm{low}\) with probability \(0.9\). A policy might choose
\[
\pi(\mathrm{wait}\mid\mathrm{low})=0.7,\qquad
\pi(\mathrm{work}\mid\mathrm{low})=0.3.
\]
The probability of the first transition \(\mathrm{low}\to\mathrm{high}\) after choosing wait is
\[
0.7\cdot0.8=0.56
\]
conditional on starting in \(\mathrm{low}\).

\paragraph{Computational neuroscience example.}
In a foraging task, the state can include patch quality, time since last reward, and internal estimates of uncertainty; actions can be stay or leave; rewards are food or juice. If the recorded state omits satiety or expectation, the Markov property may fail. Neural models often require latent state estimates, linking RL to filtering in Chapter 27 and latent variables in Chapter 34.

\paragraph{AI and machine-learning example.}
In a gridworld, states are grid locations, actions are moves, rewards are goal or collision signals, and transitions may be stochastic. In robotics, states may be positions and velocities, actions torques, and rewards task scores. In language-model RL, states can be token histories, actions next tokens, and rewards human preference scores or task outcomes.

\paragraph{Common misconceptions and failure modes.}
The reward is not the same as the value; value includes future rewards. A state is not necessarily a raw observation; it must be sufficient for prediction under the model. The Markov property does not mean independent over time. A discount factor is not only a psychological impatience parameter; it also makes infinite-horizon sums mathematically stable. A policy can be stochastic even in a deterministic environment.

\paragraph{Exercises.}
\begin{enumerate}[label=\textbf{36.1.\arabic*.}, leftmargin=*]
\item \textbf{Mechanical.} List the states, actions, transition kernel, reward variable, policy, and discount factor in a two-state MDP of your choice.
\item \textbf{Proof or derivation.} Use the chain rule and Markov assumption to derive the trajectory factorization for \(s_0,a_0,s_1,a_1,s_2\).
\item \textbf{Numerical/computational.} If rewards are \(R_1=1\), \(R_2=0\), \(R_3=2\), and \(\gamma=0.5\), compute the three-step discounted return \(1+0.5(0)+0.5^2(2)\).
\item \textbf{Interpretation.} In a neural decision task, why might a raw sensory stimulus fail to be a Markov state?
\item \textbf{Misconception/debugging.} A student says, ``The best immediate reward action is always optimal.'' Give a delayed-reward counterexample in words.
\end{enumerate}

\subsection{36.2 Value functions and Bellman equations}

\paragraph{Guiding question.}
How can future reward be summarized recursively so that long-horizon decisions become local backups?

\paragraph{Beginner-facing intuition.}
Value is expected future return. It answers: if the agent starts here and follows this policy, how much discounted reward should it expect? Bellman equations express one-step consistency. The value of a state equals expected immediate reward plus discounted value of the next state. This is the reinforcement-learning form of dynamic programming.

\paragraph{Formal definitions.}
For a policy \(\pi\), the state-value function is
\[
\boxed{
V^\pi(s)=\mathbb E_\pi[G_t\mid S_t=s].
}
\]
The action-value function is
\[
\boxed{
Q^\pi(s,a)=\mathbb E_\pi[G_t\mid S_t=s,A_t=a].
}
\]
The optimal value functions are
\[
V^*(s)=\sup_\pi V^\pi(s),\qquad Q^*(s,a)=\sup_\pi Q^\pi(s,a).
\]
For finite MDPs with bounded rewards and \(\gamma<1\), these quantities are finite.

\paragraph{Core derivation: Bellman expectation and optimality equations.}
Starting from the return identity
\[
G_t=R_{t+1}+\gamma G_{t+1},
\]
condition on \(S_t=s\) and follow policy \(\pi\):
\[
V^\pi(s)
=\mathbb E_\pi[R_{t+1}+\gamma G_{t+1}\mid S_t=s].
\]
Expanding over actions, rewards, and next states gives
\[
\boxed{
V^\pi(s)=
\sum_a\pi(a\mid s)\sum_{s',r}p(s',r\mid s,a)
\left[r+\gamma V^\pi(s')\right].
}
\]
For the optimal value, the first action should maximize expected immediate reward plus future value:
\[
\boxed{
V^*(s)=
\max_a\sum_{s',r}p(s',r\mid s,a)
\left[r+\gamma V^*(s')\right].
}
\]
Similarly,
\[
\boxed{
Q^*(s,a)=
\sum_{s',r}p(s',r\mid s,a)
\left[r+\gamma\max_{a'}Q^*(s',a')\right].
}
\]

\paragraph{Temporal-difference mechanism.}
If a transition sample \((s_t,a_t,r_{t+1},s_{t+1})\) is observed, a one-step prediction error for state values is
\[
\boxed{
\delta_t=r_{t+1}+\gamma V(s_{t+1})-V(s_t).
}
\]
TD learning updates
\[
V(s_t)\leftarrow V(s_t)+\alpha \delta_t.
\]
Q-learning uses the off-policy target
\[
Q(s_t,a_t)\leftarrow Q(s_t,a_t)+\alpha
\left[r_{t+1}+\gamma\max_{a'}Q(s_{t+1},a')-Q(s_t,a_t)\right].
\]

\paragraph{Concrete example.}
Suppose state \(s\) has one action that gives reward \(2\) and then returns to \(s\) with probability \(1\). With discount \(\gamma=0.5\),
\[
V(s)=2+0.5V(s).
\]
Thus
\[
V(s)=4.
\]
The value is larger than the immediate reward because the reward repeats.

\paragraph{Computational neuroscience example.}
In reward-learning experiments, a cue can predict later juice. A TD error \(\delta_t\) is positive when reward is better than expected and negative when reward is omitted or worse than expected. Dopamine responses in some tasks resemble reward prediction errors, but the careful statement is model-based: specific neural signals can be compared to \(\delta_t\) under a specified task model.

\paragraph{AI and machine-learning example.}
Deep Q-networks approximate \(Q(s,a)\) with a neural network when the state space is large. The Bellman target uses sampled transitions from experience replay. Actor-critic methods approximate both a policy and a value function; the critic estimates Bellman errors that guide the actor.

\paragraph{Common misconceptions and failure modes.}
Bellman equations are equations for expectations, not guarantees about one realized trajectory. Bootstrapping can propagate errors if the value estimate is poor. Q-learning's maximization can overestimate noisy action values. With function approximation, off-policy TD methods can become unstable. Rewards must be shaped carefully; optimizing the wrong reward produces the wrong behavior.

\paragraph{Exercises.}
\begin{enumerate}[label=\textbf{36.2.\arabic*.}, leftmargin=*]
\item \textbf{Mechanical.} Write the Bellman expectation equation for \(V^\pi(s)\) using \(\pi(a\mid s)\), \(p(s',r\mid s,a)\), and \(\gamma\).
\item \textbf{Proof or derivation.} Derive the Bellman equation from \(G_t=R_{t+1}+\gamma G_{t+1}\).
\item \textbf{Numerical/computational.} A state gives reward \(1\) each step and returns to itself. Compute \(V(s)\) for \(\gamma=0.9\).
\item \textbf{Interpretation.} Explain why a cue can acquire value before reward is delivered.
\item \textbf{Misconception/debugging.} A student computes TD error as \(r_{t+1}-V(s_t)\), omitting \(\gamma V(s_{t+1})\). What future-value term is missing, and why does it matter?
\end{enumerate}

\subsection{36.3 Policy optimization}

\paragraph{Guiding question.}
Instead of computing values for every action, how can we directly adjust a parameterized policy to increase expected return?

\paragraph{Beginner-facing intuition.}
Value-based methods learn how good actions are, then choose actions from those values. Policy optimization changes the policy parameters directly. This is useful when actions are continuous, when stochastic exploration is needed, or when the policy has a complex differentiable form such as a neural network.

The central trick is to differentiate an expectation over trajectories whose probabilities depend on the policy. The score-function identity turns that derivative into an expectation involving \(\nabla_\theta\log\pi_\theta(a\mid s)\).

\paragraph{Formal definitions.}
Let \(\pi_\theta(a\mid s)\) be a differentiable stochastic policy with parameters \(\theta\in\mathbb R^p\). A common objective is the expected discounted return
\[
\boxed{
J(\theta)=\mathbb E_{\tau\sim P_{\pi_\theta}}\left[\sum_{t=0}^\infty\gamma^t R_{t+1}\right].
}
\]
The trajectory distribution depends on \(\theta\) through the policy factors:
\[
P_{\pi_\theta}(\tau)=\rho(s_0)\prod_t\pi_\theta(a_t\mid s_t)P(s_{t+1}\mid s_t,a_t).
\]

\paragraph{Core derivation: score-function policy gradient.}
For a finite horizon, write \(R(\tau)\) for the total discounted return. Then
\[
J(\theta)=\sum_\tau P_{\pi_\theta}(\tau)R(\tau).
\]
Differentiate:
\[
\nabla_\theta J(\theta)
=\sum_\tau \nabla_\theta P_{\pi_\theta}(\tau)R(\tau).
\]
Use
\[
\nabla_\theta P_\theta(\tau)
=P_\theta(\tau)\nabla_\theta\log P_\theta(\tau).
\]
Since only policy terms depend on \(\theta\),
\[
\nabla_\theta\log P_\theta(\tau)
=\sum_t\nabla_\theta\log\pi_\theta(a_t\mid s_t).
\]
Therefore
\[
\boxed{
\nabla_\theta J(\theta)
=\mathbb E_{\tau\sim P_{\pi_\theta}}
\left[
R(\tau)\sum_t\nabla_\theta\log\pi_\theta(a_t\mid s_t)
\right].
}
\]
More refined policy-gradient theorems replace \(R(\tau)\) by action values or advantages to reduce variance.

\paragraph{Concrete example.}
Consider one state with two actions. Policy
\[
\pi_\theta(a=1)=\sigma(\theta),\qquad \pi_\theta(a=0)=1-\sigma(\theta)
\]
receives reward \(1\) for action \(1\) and \(0\) for action \(0\). Then
\[
J(\theta)=\sigma(\theta).
\]
The exact gradient is
\[
\nabla_\theta J=\sigma(\theta)(1-\sigma(\theta)).
\]
The score-function estimator gives reward times \(\nabla_\theta\log\pi_\theta(A)\). When \(A=1\), this is \(1-\sigma(\theta)\); when \(A=0\), the reward is \(0\). Its expectation is
\[
\sigma(\theta)(1-\sigma(\theta)),
\]
matching the exact gradient.

\paragraph{Computational neuroscience example.}
Policy-gradient ideas can model how action tendencies change when rewards are received. A softmax policy over actions resembles stochastic choice models used in behavioral neuroscience. The parameter update is a normative learning rule for a specified objective; it should not be equated directly with synaptic plasticity without a mechanistic model.

\paragraph{AI and machine-learning example.}
Policy-gradient and actor-critic methods are used in robotics, games, preference optimization, and continuous-control tasks. The actor is the policy \(\pi_\theta\). The critic estimates \(V^\pi\), \(Q^\pi\), or an advantage
\[
A^\pi(s,a)=Q^\pi(s,a)-V^\pi(s),
\]
which can reduce gradient variance and improve learning.

\paragraph{Common misconceptions and failure modes.}
Policy gradients are often high variance because rewards are delayed and trajectories are noisy. Exploration is required; actions with zero probability cannot be improved from sampled experience. A baseline can reduce variance without changing the expected gradient, but a biased critic can mislead updates. Directly maximizing sampled return can overfit to simulator quirks or reward loopholes.

\paragraph{Exercises.}
\begin{enumerate}[label=\textbf{36.3.\arabic*.}, leftmargin=*]
\item \textbf{Mechanical.} For a softmax policy \(\pi_\theta(a\mid s)\), identify the parameter, state, action, and probability being differentiated in \(\nabla_\theta\log\pi_\theta(a\mid s)\).
\item \textbf{Proof or derivation.} Prove the identity \(\nabla_\theta P_\theta(\tau)=P_\theta(\tau)\nabla_\theta\log P_\theta(\tau)\) where \(P_\theta(\tau)>0\).
\item \textbf{Numerical/computational.} In the one-state example with \(\theta=0\), compute \(J(\theta)\) and \(\nabla_\theta J(\theta)\).
\item \textbf{Interpretation.} Why can continuous-action problems favor direct policy optimization over tabular Q-learning?
\item \textbf{Misconception/debugging.} A student removes exploration after the first rewarding action is found. Explain how this can trap a policy in a suboptimal action.
\end{enumerate}

\subsection{36.4 RL and neuroscience}

\paragraph{Guiding question.}
What does reinforcement-learning mathematics clarify about neural decision-making, and where must biological interpretation be cautious?

\paragraph{Beginner-facing intuition.}
Reinforcement learning provides a compact language for reward, prediction, action, and learning from consequences. Neuroscience uses this language to model animal choice, dopamine signals, basal ganglia loops, exploration, motor learning, and cognitive control. The mathematical models are useful because they force explicit definitions: what is the state, what is the action, what is the reward, what is learned, and what information is available?

The same explicitness prevents overclaiming. A neural signal that correlates with TD error is not automatically a complete RL algorithm. A behavior that fits a model-free learner does not prove that the brain lacks a model. RL is a modeling vocabulary, not a one-line explanation of intelligence.

\paragraph{Formal definitions.}
A behavioral RL model specifies a policy
\[
\pi_\theta(a\mid s),
\]
learning dynamics for parameters or values, and an observation model for choices or neural activity. In a softmax choice model,
\[
\pi(a\mid s)=
\frac{\exp(\beta Q(s,a))}
{\sum_{a'}\exp(\beta Q(s,a'))},
\]
where \(\beta\geq0\) is an inverse-temperature parameter. Larger \(\beta\) makes choices greedier with respect to \(Q\).

A simple TD learning rule is
\[
Q(s_t,a_t)\leftarrow Q(s_t,a_t)+\alpha\delta_t,
\]
where
\[
\delta_t=r_{t+1}+\gamma\max_{a'}Q(s_{t+1},a')-Q(s_t,a_t).
\]
Model-based RL additionally learns or uses a transition model \(P(s'\mid s,a)\) and reward model, then plans with Bellman backups or search.

\paragraph{Core mechanism: reward prediction error moves value toward a target.}
For state values, the TD target is
\[
y_t=r_{t+1}+\gamma V(s_{t+1}).
\]
The error is
\[
\delta_t=y_t-V(s_t).
\]
The update
\[
V(s_t)\leftarrow V(s_t)+\alpha\delta_t
\]
is a weighted average:
\[
V_{\mathrm{new}}(s_t)=(1-\alpha)V_{\mathrm{old}}(s_t)+\alpha y_t.
\]
Thus positive prediction error raises the value estimate, negative prediction error lowers it, and \(\alpha\) controls learning speed.

\paragraph{Concrete example.}
An animal sees a cue \(C\) followed by reward \(1\). Initially \(V(C)=0\), \(V(\mathrm{after})=0\), \(\gamma=0.9\), and \(\alpha=0.2\). On the first rewarded trial,
\[
\delta=1+0.9(0)-0=1,
\]
so
\[
V(C)\leftarrow0+0.2(1)=0.2.
\]
On later trials the cue value grows, so the same reward becomes less surprising.

\paragraph{Computational neuroscience example.}
Dopamine responses in conditioning tasks often shift from unexpected reward to reward-predicting cues. TD models explain this as prediction error moving earlier in time as values are learned. Basal ganglia circuit models use action values and selection policies to describe go/no-go tendencies. Drift-diffusion and attractor models describe the neural dynamics of accumulating evidence before committing to an action. These models can be linked, but they operate at different levels: algorithmic learning rule, circuit implementation, and behavioral observation.

\paragraph{AI and machine-learning example.}
Model-free RL learns values or policies from sampled experience without explicitly planning through a learned transition model. Model-based RL learns or uses dynamics for planning, as in model predictive control. Hybrid systems combine both. Modern agents often use neural networks to represent policies, values, world models, or all three.

\paragraph{Common misconceptions and failure modes.}
Reward is a modeled scalar signal, not the same as biological pleasure. A dopamine-RPE correspondence is task- and model-dependent. Model-free and model-based are idealized categories; real agents can mix cached values, planning, habits, and inference. Neural decision variables can reflect sensory evidence, value, urgency, or motor preparation, so interpreting them requires the task model.

\paragraph{Exercises.}
\begin{enumerate}[label=\textbf{36.4.\arabic*.}, leftmargin=*]
\item \textbf{Mechanical.} In a softmax choice rule, what happens to action probabilities as \(\beta\to0\) and as \(\beta\to\infty\)?
\item \textbf{Proof or derivation.} Show that the TD update \(V\leftarrow V+\alpha(y-V)\) is a convex combination of the old value and target when \(0\leq\alpha\leq1\).
\item \textbf{Numerical/computational.} With \(V(s)=0.5\), reward \(r=1\), \(V(s')=0.4\), \(\gamma=0.9\), and \(\alpha=0.1\), compute \(\delta\) and the updated \(V(s)\).
\item \textbf{Interpretation.} Why does a reward-prediction-error model predict less neural response to an expected reward than to an unexpected reward, all else equal?
\item \textbf{Misconception/debugging.} A student says, ``Because dopamine looks like TD error, the brain is just Q-learning.'' Give two reasons this is an overclaim.
\end{enumerate}

Reinforcement learning completes the machine-learning arc from functions and distributions to action. MDPs make stochastic control explicit, Bellman equations make delayed value recursive, and policy gradients make differentiable decision rules trainable. The next volume turns the function approximators themselves into the main object: neural networks as parameterized classes of functions.

\clearpage
